% Source PDF page 23 (printed p. 9)
\section{非线性优化}
当目标函数或约束函数非线性，且不知道它们是否凸时，通常用\emph{非线性优化}（或\emph{非线性规划}）来描述这个优化问题。遗憾的是，一般非线性规划问题（1.1）没有高效的求解方法。即使问题看起来简单，只有十个变量，也可能极难求解；数百个变量的问题则可能无法处理。因此，一般非线性规划的求解方法采取了几种不同思路，每一种都需要某种妥协。

\subsection{局部优化}
在\emph{局部优化}中，我们放弃寻找在所有可行点上使目标最小的最优 $x$，转而寻找一个仅仅\emph{局部最优}的点。这意味着，它在附近的可行点中使目标函数最小，但不保证其目标值小于所有其他可行点。一般非线性规划研究中，有很大一部分集中于局部优化方法，因此这类方法已发展得相当完善。

局部优化方法可以很快，可以处理大规模问题，而且适用范围广，因为它们只要求目标函数和约束函数可微。因此，在找到一个好点即有价值、未必一定要找到最好点的应用中，局部优化方法得到了广泛使用。例如，在工程设计中，可以用局部优化改进原先由人工或其他设计方法得到的方案。

除可能找不到真正的全局最优解外，局部优化方法还有几个缺点。它们需要对优化变量给出初始猜测。这个初始猜测或起点至关重要，能够极大地影响最终局部解的目标值。算法对局部解距离全局最优还有多远，通常提供不了多少信息。局部优化方法往往对算法参数值敏感，可能需要针对某个问题或问题族调整参数。

使用局部优化方法，比求解最小二乘问题、线性规划或凸优化问题更需要技巧。使用者需要尝试选择算法、调整算法参数，并在求解单个实例时找到足够好的初始猜测，或者在求解一族问题时找到产生足够好初始猜测的方法。大致来说，局部优化方法更像一门技艺，而非成熟技术。它是一门发展完善、往往十分有效的技艺，但终究仍是技艺。相比之下，最小二乘和线性规划的求解几乎不需要这种技巧，当然，接近当前计算能力极限的问题除外。

将非线性规划的局部优化方法与凸优化比较，会发现一个有趣的区别。大多数局部优化方法只要求目标函数和约束函数可微，因此将实际问题写成非线性优化问题相对直接。局部优化的技巧，在于建模之后如何求解，或者说如何找到局部最优点。

% Source PDF page 24 (printed p. 10)
凸优化恰好相反：技巧和挑战在于问题建模；一旦将问题写成凸优化问题，求解便相对直接。

\subsection{全局优化}
\emph{全局优化}寻找优化问题（1.1）真正的全局解，代价是效率。全局优化方法的最坏情形复杂度随问题规模 $n$ 和 $m$ 指数增长；人们希望，对于实际遇到的特定问题实例，方法的运行速度能快得多。这种有利情形确实存在，但并不常见。即使只有几十个变量的小问题，也可能需要很长时间，例如数小时或数天，才能求解。

全局优化用于变量少、计算时间不紧迫，而且找到真正全局解具有很高价值的问题。工程设计中的一个例子，是高价值系统或安全关键系统的最坏情形分析与验证。这里的变量表示在制造过程中，或随环境、运行条件而变化的不确定参数。目标函数是效用函数，即数值越小表示性能越差；约束表示关于参数可能取值的先验知识。优化问题（1.1）就是寻找最坏的参数值。若最坏情况下的性能仍然可以接受，就能证明系统对于这些参数变化是安全或可靠的。

局部优化方法可以迅速找到一组较差的参数，但不能保证它们是绝对最坏的。如果局部优化找到使性能不可接受的参数，就成功判定了系统不可靠。但局部优化无法证明系统可靠；它可能只是没有找到坏的参数值。相比之下，全局优化方法会找到绝对最坏的参数值，如果相应性能可以接受，就能证明系统安全。代价是计算时间，即使参数相对较少，计算时间也可能很长。不过，当性能认证本身的价值很高，或对可靠性、安全性判断错误的代价很高时，这种计算仍可能值得。

\subsection{凸优化在非凸问题中的作用}
本书主要讨论凸优化问题，以及能归结为凸优化的应用。但凸优化在非凸问题中也起着重要作用。

% Source PDF page 25 (printed p. 11)
\subsubsection*{为局部优化提供初值}
一种显然的用法是将凸优化与局部优化结合。从一个非凸问题出发，先找出一个近似的凸模型。这个近似问题容易求解，也不需要初始猜测；求解后，便得到近似凸问题的精确解。再以此点为起点，对原来的非凸问题应用局部优化方法。

\subsubsection*{非凸优化的凸启发式方法}
若干求解非凸问题的启发式方法以凸优化为基础。我们将看到一个有趣的例子：寻找满足某些约束的稀疏向量 $x$，即非零分量很少的向量。尽管这是一个困难的组合问题，一些基于凸优化的简单启发式方法往往能找到相当稀疏的解。（这些方法见第6章。）

另一大类例子是随机化算法：从某个概率分布中抽取若干候选点，再把其中最好的点作为非凸问题的近似解。假设用于抽取候选点的分布族由某些参数刻画，例如均值和协方差。于是可以提出问题：哪一个分布能使目标函数的期望值最小？这个问题有时是凸问题，因此能够高效求解。（例如参见习题11.23。）

\subsubsection*{全局优化中的界}
许多全局优化方法需要以很小的计算代价，得到非凸问题最优值的下界。两种标准方法都以凸优化为基础。在\emph{松弛}中，用一个更宽松的凸约束替代每个非凸约束。在\emph{拉格朗日松弛}中，则求解第5章介绍的拉格朗日对偶问题。对偶问题是凸问题，并给出非凸问题最优值的下界。

\section{全书结构}
本书分为三个主要部分，分别是理论、应用和算法。

\subsection{第一部分：理论}
第一部分“理论”介绍凸分析和凸优化的基本定义、概念与结果。我们不追求百科全书式的覆盖，而是着重选择那些有助于识别和建立凸优化问题的内容。这些都是经典材料，几乎全部可以在其他凸分析和优化教材中找到。

% Source PDF page 26 (printed p. 12)
我们也不追求给出最一般形式的结果；对此，读者可以参考任一标准凸分析教材。

第2章和第3章分别讨论凸集与凸函数。除了给出一些常见例子，还介绍若干凸性运算规则，即保持集合或函数凸性的运算。把基本例子与这些规则结合起来，就能构造，或者更重要地，识别一些相当复杂的凸集与凸函数。

第4章“凸优化问题”严格讨论优化问题，并介绍若干可用于重写问题的变换。我们还介绍一些常见的凸优化子类，例如线性规划、几何规划，以及较新发展的二阶锥规划和半正定规划。

第5章讨论拉格朗日对偶，它在凸优化中起着核心作用。这里给出经典的 Karush--Kuhn--Tucker 最优性条件，以及凸优化问题的局部与全局敏感性分析。

\subsection{第二部分：应用}
第二部分“应用”介绍概率与统计、计算几何、数据拟合等领域中的多种凸优化应用。我们希望，对这些应用的叙述能为广泛的读者所理解。为了保持每个应用简短，我们只考虑简单情形，有时补充说明可能的扩展。我们确信，某些应用的处理方式会令专家不以为然，在此先向他们致歉。但我们的目的是迅速地向广泛读者传达这些应用的基本面貌，而非作优雅、理论严密或完整的论述。我们自己的背景是电气工程中的控制系统、信号处理、电路分析与设计。尽管在以本书为主要教材的课程中，我们也讲授这些内容，但其中只有少数应用具有足够广泛的可理解性，因而被收入本书。

第二部分的目标，是通过例子向读者展示如何在实际中应用凸优化。

\subsection{第三部分：算法}
第三部分“算法”介绍求解凸优化问题的数值方法，重点是牛顿法与内点法。它包括三章，依次讨论无约束优化、等式约束优化和不等式约束优化。这些章节遵循一个自然的层次结构：把一个问题的求解，归结为求解一系列更简单的问题。二次优化问题，包括最小二乘等，构成这一层次的基础；它们可以通过求解一组线性方程而精确求解。

% Source PDF page 27 (printed p. 13)
第9章与第10章介绍的牛顿法，是这一层次的下一级。牛顿法将无约束或等式约束问题的求解，归结为求解一系列二次问题。第11章介绍位于最高层的内点法。这类方法通过求解一系列无约束或等式约束问题，来求解不等式约束问题。

总体而言，我们只讨论少数几个算法，省略了拟牛顿、共轭梯度、束方法和割平面等整类优秀方法。即使是所介绍的方法，我们也采用简化变体，而非最新、最复杂的版本。算法的选择遵循几个标准：描述与实现简单，同时可靠、稳健，并且对大多数问题足够有效、足够快。

许多凸优化使用者最终会使用标准软件，例如线性或半正定规划求解器，而不会亲自开发这类软件。对这些读者，第三部分旨在传达各种方法的基本面貌及其主要特征。对少数最终会开发新算法的读者，我们认为第三部分也能作为良好的入门。

\subsection{附录}
本书有三个附录。第一个列出正文中使用的一些基本数学事实，同时也用于说明记号。第二个讨论一个相当专门的主题：目标为二次函数、且只有一个二次约束的优化问题。这些问题虽然非凸，却能高效求解；第二部分介绍的若干应用会用到这些结果。

最后一个附录简要介绍数值线性代数，重点讨论能够利用稀疏性等问题结构提高效率的方法。我们没有涵盖若干重要主题，例如舍入误差分析，也没有详细介绍所需矩阵分解的实现方法。这些主题已有多本优秀教材讨论。

\subsection{关于例子的说明}
正文许多地方，尤其介绍应用与算法的第二、三部分，使用具体例子说明思想。有些例子专门选择或设计来说明某个观点；另一些则选择为“典型”例子，也就是从某个显然或简单的概率分布中抽取的样本。凭几十个或几百个随机生成的例子推断算法性能的风险，已广为人知，这里不再重复。这些例子只是粗略展示算法性能，或计算量随问题维数变化的大致情况，并非用于准确预测算法性能。尤其是，读者得到的结果可能与我们的不同。

% Source PDF page 28 (printed p. 14)
\subsection{关于习题的说明}
每章末尾都有一组习题。有些要求补全正文论证或断言的细节；另一些要求判断或证明给定集合、函数、问题的凸性，或者更一般地，建立凸优化问题。某些章节含有数值习题，需要使用合适的高级语言进行少量编程。习题难度混合，不事先说明，既有相当直接的问题，也有颇具技巧的问题。

\section{记号}
除少数例外，本书记号基本标准。本节介绍基本记号；第697页列出更完整的记号表。

我们用 $\R$ 表示实数集，$\R_+$ 表示非负实数集，$\R_{++}$ 表示正实数集。实 $n$ 维向量的集合记作 $\R^n$，实 $m\times n$ 矩阵的集合记作 $\R^{m\times n}$。向量和矩阵用方括号括起，分量之间用空格隔开。用圆括号括起逗号分隔的元素列表，表示由它们构成的列向量。例如，若 $a,b,c\in\R$，则
\[
 (a,b,c)=\begin{bmatrix}a\\b\\c\end{bmatrix}
 =\begin{bmatrix}a&b&c\end{bmatrix}^{\T}\in\R^3.
\]
符号 $\one$ 表示各分量均为一的向量，维数由上下文确定。记号 $x_i$ 既可表示向量 $x$ 的第 $i$ 个分量，也可表示向量集合或序列 $x_1,x_2,\ldots$ 中的第 $i$ 个元素，具体含义由上下文或正文说明。

用 $\mathbf S^k$ 表示对称 $k\times k$ 矩阵集，$\mathbf S_+^k$ 表示对称半正定 $k\times k$ 矩阵集，$\mathbf S_{++}^k$ 表示对称正定 $k\times k$ 矩阵集。弯曲不等号 $\succeq$（以及严格形式 $\succ$）表示\emph{广义不等式}：对向量表示逐分量不等式，对对称矩阵则表示矩阵不等式。带下标的 $\preceq_K$（或 $\prec_K$）表示相对于锥 $K$ 的广义不等式，见第2.4.1节。

我们描述函数的记号与通常约定略有不同，但希望不会引起混淆。记号 $f:\R^p\to\R^q$ 表示：$f$ 在 $\R^p$ 的某个子集上取值于 $\R^q$，这个子集就是它的定义域，记作 $\dom f$。可以像计算机语言中的类型声明那样理解这个记号：$f:\R^p\to\R^q$ 表示函数 $f$ 接受一个实 $p$ 维向量作为输入，返回一个实 $q$ 维向量。定义域 $\dom f$ 指明 $\R^p$ 中哪些点 $x$ 的 $f(x)$ 有定义。例如，我们把对数函数写成 $\log:\R\to\R$，且 $\dom\log=\R_{++}$。

% Source PDF page 29 (printed p. 15)
记号 $\log:\R\to\R$ 表示对数函数接受并返回实数；$\dom\log=\R_{++}$ 表示它仅对正数有定义。

我们用 $\R^n$ 代表一般的有限维向量空间。书中还会遇到其他有限维向量空间，例如次数不超过给定上限的一元多项式空间，以及对称 $k\times k$ 矩阵空间 $\mathbf S^k$。选定向量空间的一组基，就总能把它与 $\R^n$ 对应起来，其中 $n$ 是其维数，因此对 $\R^n$ 叙述的一般结果也能应用于这些空间。我们通常将一般结论向其他向量空间的转换留给读者。例如，任意线性函数 $f:\R^n\to\R$ 都可以写成 $f(x)=c^{\T}x$，其中 $c\in\R^n$。对 $\mathbf S^k$，选定一组基并作相应转换，就得到：任意线性函数 $f:\mathbf S^k\to\R$ 都可以写成 $f(X)=\tr(CX)$，其中 $C\in\mathbf S^k$。

% Source PDF page 30 (printed p. 16)
\section*{文献说明}
\addcontentsline{toc}{section}{文献说明}
最小二乘是一个非常古老的主题，例如 Gauss 在19世纪20年代用拉丁文写成的专著，近期由 Stewart 译出 [Gau95]。较新的著作包括 Lawson 和 Hanson [LH95] 以及 Bj\"orck [Bj\"o96] 的书。线性规划的参考文献见第4章。

非线性规划的局部方法有许多优秀教材，包括 Gill、Murray 和 Wright [GMW81]，Nocedal 和 Wright [NW99]，Luenberger [Lue84]，以及 Bertsekas [Ber99]。

全局优化可参阅 Horst 和 Pardalos [HP94]、Pinter [Pin95]、Tuy [Tuy98] 的著作。用凸优化为非凸问题求界，是一个活跃的研究主题；上述全局优化著作、Ben-Tal 和 Nemirovski 的书 [BTN01，第4.3节]，以及 Nesterov、Wolkowicz 和 Ye 的综述 [NWY00] 都有讨论。相关的重要论文包括 Goemans 和 Williamson [GW95]、Nesterov [Nes00, Nes98]、Ye [Ye99] 和 Parrilo [Par03]。随机化方法见 Motwani 和 Raghavan [MR95]。

凸分析研究凸集、凸函数和凸优化问题的数学理论，是一个发展完善的数学分支。基本参考书包括 Rockafellar [Roc70]、Hiriart-Urruty 和 Lemar\'echal [HUL93, HUL01]、Borwein 和 Lewis [BL00]，以及 Bertsekas、Nedi\'c 和 Ozdaglar [Ber03]。更多凸分析文献见第2至第5章。

Nesterov 和 Nemirovski [NN94] 首先指出，内点法能够求解许多凸优化问题；另见第11章的参考文献。Ben-Tal 和 Nemirovski 的书 [BTN01] 讨论现代凸优化、内点法及其应用。

本书未介绍的凸优化求解方法包括次梯度法 [Sho85]、束方法 [HUL93]、割平面法 [Kel60, EM75, GLY96] 和椭球法 [Sho91, BGT81]。

凸优化问题易于处理的思想并不新鲜。人们早已认识到，凸优化理论比一般非线性优化理论直接得多，也完整得多。正是在这一背景下，Rockafellar 在1993年发表的 SIAM Review 综述 [Roc93] 中写道：
\begin{quote}
事实上，优化的重大分水岭不在线性与非线性之间，而在凸性与非凸性之间。
\end{quote}

Nemirovski 和 Yudin 于1983年在其著作 \emph{Problem Complexity and Method Efficiency in Optimization} [NY83] 中，首次正式论证了凸优化问题比一般非线性优化问题更容易求解。他们证明，凸优化问题的基于信息的复杂度远低于一般非线性优化问题。较新的相关著作为 Vavasis [Vav91]。

内点法较低的理论复杂度，是这一领域现代研究不可分割的一部分。许多研究致力于证明：某种内点法或其他方法，求解某类凸优化问题所需的运算次数，随问题维数和 $\log(1/\epsilon)$ 的增长至多为多项式，其中 $\epsilon>0$ 是所需精度。（第11章将介绍若干简单结果。）Nesterov 和 Nemirovski 的书 [NN94] 首次对此作了全面讨论。其他著作包括 Ben-Tal 和 Nemirovski [BTN01，第5讲] 和 Renegar [Ren01]。内点法对各种凸优化问题的多项式时间复杂度，与许多非凸优化问题形成鲜明对比：对于后者，所有已知算法在最坏情况下所需的运算次数，都随问题维数指数增长。

% Source PDF page 31 (printed p. 17)
凸优化已用于众多领域，无法在此逐一列举。凸分析在经济与金融中居于核心地位，是许多结果的基础。例如，将分离超平面定理与无套利假设结合，可以推出价格与风险中性概率的存在性，参见 Luenberger [Lue95, Lue98] 和 Ross [Ros99]。近期，凸优化，尤其是半正定规划的可求解性，在自动控制理论中受到了特别关注。控制理论中的凸优化应用，可参阅 Boyd 和 Barratt [BB91]，Boyd、El Ghaoui、Feron 和 Balakrishnan [BEFB94]，Dahleh 和 Diaz-Bobillo [DDB95]，El Ghaoui 和 Niculescu [EN00]，以及 Dullerud 和 Paganini [DP00]。嵌入式凸优化的一个良好例子是模型预测控制；这项自动控制技术要求每一步求解一个凸二次规划。模型预测控制现已广泛用于化工过程控制行业，见 Morari 和 Zafiriou [MZ89]。另一个具有悠久凸优化，尤其是几何规划应用历史的领域，是电子电路设计。相关论文包括 Fishburn 和 Dunlop [FD85]，Sapatnekar、Rao、Vaidya 和 Kang [SRVK93]，以及 Hershenson、Boyd 和 Lee [HBL01]。Luo [Luo03] 综述了信号处理与通信中的应用。更多凸优化应用文献见第4章及第6至第8章。

近期凸优化内点法的高质量实现，可在 LOQO [Van97]、MOSEK [MOS02] 软件包及第11章列出的程序中找到。用于描述优化问题的软件系统包括 AMPL [FGK99] 和 GAMS [BKMR98]。二者都在一定程度上支持识别能够转化为线性规划的问题。

% Source PDF page 32 (printed p. 18; blank source page)
